% from ega3-1-fr.tex
\oldpage[III]{117}
En effet, soit
$0\to\mathcal F'\to\mathcal F\to\mathcal F''\to0$ une suite exacte de
$\mathcal O_X$-Modules cohérents~; supposons par exemple que $\mathcal F'$
et $\mathcal F''$ appartiennent à $K'$~; on a la suite exacte de cohomologie
\[
R^{q-1}f_*(\mathcal F'')\xrightarrow{\partial}R^qf_*(\mathcal F')
\longrightarrow R^qf_*(\mathcal F)\longrightarrow R^qf_*(\mathcal F'')
\xrightarrow{\partial}R^{q+1}f_*(\mathcal F')
\]
dans laquelle par hypothèse les quatre termes extrêmes sont cohérents~; il
en est donc de même du terme médian $R^qf_*(\mathcal F)$ par
(0\textsubscript{I}, 5.3.4 et 5.3.3). On montre de la même manière que
lorsque $\mathcal F$ et $\mathcal F'$ (resp. $\mathcal F$ et
$\mathcal F''$) sont dans $K'$, il en est de même de $\mathcal F''$ (resp.
$\mathcal F'$). De plus, tout \emph{facteur direct} cohérent $\mathcal F'$
d'un $\mathcal O_X$-Module $\mathcal F\in K'$ appartient aussi à $K'$~: en
effet, $R^qf_*(\mathcal F')$ est alors facteur direct de
$R^qf_*(\mathcal F)$ (G, II, 4.4.4), donc est de type fini, et comme il est
quasi-cohérent (\hyperref[III.1.4.10-fr]{1.4.10}), il est cohérent, $Y$
étant noethérien. En vertu de (\hyperref[III.3.1.3-fr]{3.1.3}), on est
ramené à prouver que lorsque $X$ est \emph{irréductible} de point générique
$x$, il existe un $\mathcal O_X$-Module cohérent $\mathcal F$ appartenant à
$K'$, tel que $\mathcal F_x\neq0$~: en effet, si ce point est établi, on
pourra l'appliquer à tout sous-préschéma fermé irréductible $Y$ de $X$, car
si $j:Y\to X$ est l'injection canonique, $f\circ j$ est propre (II, 5.4.2),
et si $\mathcal G$ est un $\mathcal O_Y$-Module cohérent de support $Y$,
$j_*(\mathcal G)$ est un $\mathcal O_X$-Module cohérent tel que
\[
R^q(f\circ j)_*(\mathcal G)=R^qf_*(j_*(\mathcal G))
\]
(G, II, 4.9.1), donc on est bien dans les conditions d'application de
(\hyperref[III.3.1.3-fr]{3.1.3}).

Or, en vertu du lemme de Chow (II, 5.6.2), il existe un préschéma
irréductible $X'$ et un morphisme \emph{projectif} et surjectif
$g:X'\to X$ tel que $f\circ g:X'\to Y$ soit \emph{projectif}. Il existe un
$\mathcal O_{X'}$-Module $\mathcal L$ ample pour $g$ (II, 5.3.1)~;
appliquons le th. fondamental des morphismes projectifs
(\hyperref[III.2.2.1-fr]{2.2.1}) à $g:X'\to X$ et à $\mathcal L$~: il
existe donc un entier $n$ tel que
$\mathcal F=g_*(\mathcal O_{X'}(n))$ soit un $\mathcal O_X$-Module cohérent
et $R^qg_*(\mathcal O_{X'}(n))=0$ pour tout $q>0$~; en outre, comme
\[
g^*(g_*(\mathcal O_{X'}(n)))\longrightarrow\mathcal O_{X'}(n)
\]
est surjectif pour $n$ assez grand (\hyperref[III.2.2.1-fr]{2.2.1}), on
voit qu'on peut supposer que, au point générique $x$ de $X$, on a
$\mathcal F_x\neq0$ (II, 3.4.7). D'autre part, comme $f\circ g$ est
projectif et $Y$ noethérien, les
$R^p(f\circ g)_*(\mathcal O_{X'}(n))$ sont \emph{cohérents}
(\hyperref[III.2.2.1-fr]{2.2.1}). Cela étant,
$R^*(f\circ g)_*(\mathcal O_{X'}(n))$ est l'aboutissement d'une suite
spectrale de Leray, dont le terme $E_2$ est donné par
\[
E_2^{pq}=R^pf_*(R^qg_*(\mathcal O_{X'}(n)))~;
\]
ce qui précède montre que cette suite spectrale est dégénérée, et on sait
alors (0, 11.1.6) que $E_2^{p0}=R^pf_*(\mathcal F)$ est isomorphe à
$R^p(f\circ g)_*(\mathcal O_{X'}(n))$, ce qui achève la démonstration.
\end{proof}

\begin{corollary}[3.2.2]
\label{III.3.2.2-fr}
Soit $Y$ un préschéma localement noethérien. Pour tout morphisme propre
$f:X\to Y$, l'image directe par $f$ de tout $\mathcal O_X$-Module cohérent
est un $\mathcal O_Y$-Module cohérent.
\end{corollary}

\begin{corollary}[3.2.3]
\label{III.3.2.3-fr}
Soient $A$ un anneau noethérien, $X$ un schéma propre sur $A$~; pour tout
$\mathcal O_X$-Module cohérent $\mathcal F$, les $H^p(X,\mathcal F)$ sont
des $A$-modules de type fini, et il existe un entier $r>0$ tel que pour tout
$\mathcal O_X$-Module cohérent $\mathcal F$ et tout $p>r$,
$H^p(X,\mathcal F)=0$.
\end{corollary}

\begin{proof}
La seconde assertion a déjà été démontrée
(\hyperref[III.1.4.12-fr]{1.4.12})~; la première résulte du th. de finitude
(\hyperref[III.3.2.1-fr]{3.2.1}), compte tenu de
(\hyperref[III.1.4.11-fr]{1.4.11}).
\end{proof}

En particulier, si $X$ est un \emph{schéma algébrique propre} sur un corps
$k$, alors, pour tout $\mathcal O_X$-Module cohérent $\mathcal F$, les
$H^p(X,\mathcal F)$ sont des $k$-espaces vectoriels de \emph{dimension
finie}.

\begin{corollary}[3.2.4]
\label{III.3.2.4-fr}
Soient $Y$ un préschéma localement noethérien, $f:X\to Y$ un morphisme de
type fini. Pour tout $\mathcal O_X$-Module cohérent $\mathcal F$ dont le
support est propre sur $Y$ (II, 5.4.10), les $\mathcal O_Y$-Modules
$R^qf_*(\mathcal F)$ sont cohérents.
\end{corollary}

